\documentclass{article}
\usepackage{amsmath}

\title{A Novel Framework for Graph Optimization}
\author{Test Author}
\date{}

\begin{document}
\maketitle

\section{Introduction}
Graphs are used in many applications.
Many algorithms have been proposed.
Our algorithm is better.
We test it on several datasets.

\section{Related Work}
Smith et al.\ (2020) proposed a method for graph partitioning.
Jones et al.\ (2021) extended this to weighted graphs.
There are also methods based on spectral clustering.

\section{Our Method}
We propose a new framework for graph optimization.
The framework has three phases.

In Phase 1, we compute a preliminary partition using a greedy approach.
In Phase 3, we refine the partition using local search.
In Phase 2, we merge small clusters.

The time complexity is $O(n \log n)$ in the best case.
However, in our experiments the runtime was often quadratic.

\section{Experiments}
We tested our method on five benchmark datasets.
Our method achieved the best results on three of them.
On the other two datasets, baseline methods performed better.
We attribute this to the nature of the graphs, but we did not investigate further.

Table~1 shows the results but we were unable to include it due to space constraints.

\section{Theoretical Analysis}

\begin{theorem}
The proposed method converges to a local optimum in polynomial time.
\end{theorem}

We do not prove this theorem formally but it seems intuitively obvious from the structure of the algorithm.

\section{Conclusion}
We presented a novel framework for graph optimization that outperforms existing methods in most cases.
Future work will address the limitations identified in our experiments.
Our contribution is significant because graph optimization is important.

\end{document}
